\section{Introduction}
\label{sec:introduction}

A mean-payoff parity game asks a player to satisfy both a long-run
reward constraint and a recurring-priority condition. These objectives
interact: repeatedly visiting a favorable priority may incur a loss,
even when a nonnegative mean reward is achievable on its own. We show
that the combined winning region can be computed in deterministic
quasipolynomial time in the full binary input length.

An instance consists of a finite directed graph $G=(V,E)$ in which
every vertex has an outgoing edge, a partition
$V=V_{\Max}\sqcup V_{\Min}$, a reward $r(v)\in\Z$ at each vertex,
and a priority $p(v)\in\N$ at each vertex, where $0\in\N$.
The graph with its ownership partition is called an \emph{arena}.
The owner of the current vertex chooses the next edge. Strategies are
pure: each chooses an outgoing edge after every finite history ending
at a vertex owned by that player. A strategy is \emph{positional} if
the choice depends only on the current vertex. Max wins an infinite
play $\pi=v_0v_1\cdots$ precisely when
\begin{equation}\label{eq:objective}
 \min\{p(v):v\text{ occurs infinitely often in }\pi\}
 \text{ is even},\qquad
 \MP(\pi):=\liminf_{T\to\infty}
       \frac1T\sum_{t=0}^{T-1}r(v_t)\ge0.
\end{equation}
Min's winning objective is the complement of this conjunction.
Let $L$ be the length of a fixed conventional explicit binary encoding
of the graph, ownership, rewards, and priorities. In particular, large
reward magnitudes and large priority values are charged by their
binary lengths.

\begin{theorem}\label{thm:main}
A uniform deterministic algorithm returns exactly all vertices from
which Max can force~\eqref{eq:objective} against every Min strategy.
For an absolute constant $C>0$, it uses at most
\[
                     2^{C(\log_2(L+2))^2}
\]
bit operations. The number of distinct priorities and the magnitudes
of the rewards are unrestricted.
\end{theorem}

The algorithmic input is the deterministic mean-payoff theorem of
\cite[Theorem~1.1]{OpenAI2026Deterministic}, together with its
positional-strategy guarantees
\cite[Corollary~7.1]{OpenAI2026Deterministic}.
We state the exact interface in Proposition~\ref{prop:mp-interface}.
The reduction from the conjunction to ordinary mean payoff is proved
here, including the determinacy and strategy-combination facts used
by the recursion. Its number of mean-payoff queries is
$2^{O((\log(n+2))^2)}$, where $n=|V|$, and every queried arena has
encoding length polynomial in $L+2$.

\paragraph{The reduction.}
A \emph{dominion} for a player is a set of vertices within which that
player can keep play and win from every start. Following the
reduced-precision approach to parity games~\cite{Parys2019}, we compute tentative
labels that must be correct on every dominion below a specified size
bound for each player. Outside those protected dominions, the labels
may be incorrect. A call has at most one recursive child retaining both
bounds; every other child halves one bound. Each child removes the
parent's least priority. These two measures give the quasipolynomial
call count.

The essential additional issue is how to combine strategies without
spoiling mean payoff. A limiting average bound for each uninterrupted
play does not itself control the cost of arbitrarily many restarts.
We instead construct Max strategies whose every consistent prefix of
length $t$ has weight at least $-\epsilon t-C_\epsilon$, for every
$\epsilon>0$, with $C_\epsilon$ uniform over the opponent's choices.
Using the increasing-round construction
of~\cite[Lemma~7 of the full version]{BouyerMarkeyOlschewskiUmmels2011},
between attempts to visit an even least priority the strategy follows
a mean-payoff policy for $1,2,3,\ldots$ edges. The number of such rounds
completed by time $T$ is $O(\sqrt T)$, so bounded overhead per round
has vanishing average. This uniform estimate supports both
determinacy and the dominion extraction needed by the recursion.

The two ingredients are separated in the proof: Section~\ref{sec:strategies}
establishes the strategy and dominion statements;
Section~\ref{sec:algorithm} gives the bounded-dominion algorithm;
and Section~\ref{sec:complexity} counts its queries and bit operations.
The recursion uses the ordinary mean-payoff solver only through its
winning-set interface, allowing that solver to be replaced without
changing the parity part of the construction.

\subsection{Historical context and the role of this reduction}
\label{sec:history}

Chatterjee, Thomas A.~Henzinger, and Jurdzi\'nski introduced
mean-payoff parity games to combine qualitative specifications with
quantitative performance requirements
\cite{ChatterjeeHenzingerJurdzinski2005}. They established the existence
of optimal strategies and showed that the player enforcing the
conjunction may need infinite memory. Bouyer, Markey, Olschewski,
and Ummels proved that the opposing player has an optimal positional
strategy and gave a simpler recursive value algorithm
\cite[Theorems~2 and~7]{BouyerMarkeyOlschewskiUmmels2011}.
Thus the strategy asymmetry and the recursive decomposition of these
games precede the present work.

Chatterjee and Doyen related mean-payoff parity games to energy parity
games, whose accumulated weight must remain above a fixed lower bound,
and established membership of the threshold problem in
$\mathrm{NP}\cap\mathrm{coNP}$~\cite{ChatterjeeDoyen2012}.
Chatterjee, Monika Henzinger, and Svozil improved the threshold
running time to $O(n^{d-1}mW)$, for $n$ vertices, $m$ edges,
$d\ge2$ priority levels, and $W$ the maximum of $1$ and the absolute weights
\cite[Theorem~6]{ChatterjeeHenzingerSvozil2017}.
This dependence is polynomial in the numerical weights but
exponential in the number of priorities. Daviaud, Jurdzi\'nski,
and Lazi\'c replaced the latter dependence by a quasipolynomial
factor, using strategy decompositions and progress measures combining
parity and energy information~\cite{DaviaudJurdzinskiLazic2018}.
Anand, Fijalkow, Goubault-Larrecq, Leroux, and Ohlmann subsequently
gave a separating-automaton construction for a related disjunction
of parity and mean-payoff objectives
with comparable pseudo-quasipolynomial complexity
\cite{AnandEtAl2021}. These bounds retain a numerical weight factor.

For parity games alone, Calude, Jain, Khoussainov, Li, and Stephan
obtained a quasipolynomial algorithm
\cite{CaludeEtAl2017}; Jurdzi\'nski and Lazi\'c developed the
succinct progress measures underlying the later mean-payoff parity
algorithm~\cite{JurdzinskiLazic2017}. The recursive ancestry relevant
here is Parys's modification of Zielonka's attractor recursion
\cite{Parys2019}, developed further by Lehtinen, Parys, Schewe,
and Wojtczak~\cite{LehtinenParysScheweWojtczak2022}.
Parys's Algorithm~2 and Lemma~6.1 already use separate size bounds
for the two players' dominions, two sequences of reduced-precision
calls surrounding one full-precision call, and the associated
dominion-preservation argument. Our algorithm adapts this structure
by repeatedly removing vertices where Max loses the ordinary
mean-payoff test. The strategy lemmas justify that adaptation for
the conjunction.

The strategy construction also has direct predecessors.
Increasing payoff phases interspersed with parity visits occur in
\cite{ChatterjeeHenzingerJurdzinski2005}; the even-priority combination
implication is stated in
\cite[Lemma~5]{BouyerMarkeyOlschewskiUmmels2011}.
The full version of that paper proves it with a phase of length $j$
in round $j$. The same scheduling appears in the threshold recursion
of~\cite[Appendix~C, full version]{ChatterjeeHenzingerSvozil2017}.
Daviaud, Jurdzi\'nski, and Lazi\'c also maintain a uniform bound on
finite-prefix costs in their strategy-decomposition proof, for their
strict mean-payoff convention
\cite[Sections~2.3--2.4]{DaviaudJurdzinskiLazic2018}.
Here we make the all-$\epsilon$ prefix guarantee explicit for a
single strategy at the non-strict threshold and use it to control
repeated restarts in the increasing-round construction against
history-dependent opponents.

The resulting reduction uses quasipolynomially many ordinary
mean-payoff queries of polynomial binary length. Substituting the
deterministic mean-payoff algorithm
\cite[Theorem~1.1 and Corollary~7.1]{OpenAI2026Deterministic}
therefore gives the full binary-length bound. This consequence uses
the present query bound together with that companion theorem; the
earlier recursive, strategy-combination, and dominion arguments supply
its methodological foundations.

